\section*{Resumo}
\addcontentsline{toc}{section}{Resumo}

\textbf{Contexto.} PREENCHER: origem da série (extrato, síntese), alvo terapêutico e dado experimental que motiva o estudo (ex.: \% de inibição ou $\text{IC}_{50}$ de cada composto).

\textbf{Objetivo.} PREENCHER: a pergunta central, em uma frase (ex.: que modo de ligação explica a atividade do composto líder frente aos análogos?).

\textbf{Metodologia.} PREENCHER: estruturas cristalográficas (código PDB), microespécies no pH do ensaio, amostragem por difusão (DiffDock-L), minimização local (GNINA), validação física (PoseBusters), seleção da pose no sítio ativo e mapeamento de interações (PROPKA3 + ProLIF).

\textbf{Resultados e discussão.} PREENCHER: RMSD do \textit{redocking}; métrica primária pré-registrada e as demais métricas; modo de ligação do líder e o que o distingue dos análogos; interações-chave atribuídas ao fragmento do ligante.

\textbf{Conclusões.} PREENCHER: hipóteses estruturais testáveis e as limitações do modelo (receptor rígido, ausência de águas explícitas, escores semiquantitativos, microespécie).

\vspace{6pt}
\textbf{Palavras-chave:} PREENCHER; \textit{docking} molecular; relação estrutura--atividade.
